\section{Limitations}\label{sec:limits}
\begin{itemize}
\item \textbf{Synthetic scope.} Both benchmarks, including the rendered one, are synthetic, low-dimensional and
compact, with known $K$ and orthogonal actions. We make no claim about natural images, physical systems or large models.
\item \textbf{No identifiability.} Wrong closed algebras are frequent for SO(3), and correct algebras still have
irreducible action ambiguity (Prop.~\ref{prop:endpoint}c). Our ground-truth comparisons use orthogonal Procrustes, so
they assume a near-isometric encoder; the $T^2$ failures are seed-level fitting failures shared by all methods.
\item \textbf{Limited baselines.} We compare controlled objective variants and an action-supervised oracle; we do not
reimplement published methods with different supervision.
\item \textbf{Statistics and protocol.} We use 20 seeds; effects outside SO(3) are small or concentrated in few
seeds. No transport improvement is established; on rendered images transport is slightly worse. The closure-weight rule needed an amendment, committed before any fresh
run. One pre-registered prediction (structure beats CKA at predicting transfer failure) was not supported.
\item \textbf{Theory is local, qualitative and self-checked.} The escape bound covers only small actions; the
classification is specific to $\so(4)$ and $\so(5)$; Prop.~\ref{prop:endpoint} concerns exact zeros at the
ground-truth embedding, gives no quantitative gap and says nothing about the optimization path. The proofs have not
been checked independently of the AI system that drafted them.
\end{itemize}

\section{Conclusion}\label{sec:conclusion}
Closure turns continuous disagreement between transformation models into a choice among a few closed subalgebras.
For SO(3) the choice is often wrong, because endpoints cannot reveal rotations about the axis through a point.
Agreement, of embeddings or of transformations, is not evidence of correctness.
